\section{Discussion}
\label{sec:discussion}

Our experiment refutes the simple form of the hypothesis --- graded, distance-monotonic
robustness with orthogonal basins --- and in doing so surfaces two findings the original
framing did not anticipate.

\para{Orthogonal is not robust.} Far-from-domain capability prompts drifted the \emph{most},
not the least. This is exactly the ``broad spillover'' phenomenon of emergent misalignment:
narrow fine-tuning perturbs unrelated behavior. It aligns with \citet{springer2026geometry},
who argue orthogonality to the fine-tuning objective confers no stability, and against the
naive reading of the safety basin as a behavioral attractor on orthogonal tasks. The real
low-drift region in our data is organized by \emph{content type} --- factual and technical
knowledge --- rather than by \emph{task distance}. A model's stored facts about SQL injection
or password storage survive fine-tuning even though those questions are semantically far from
the code-completion training data, while a math word problem at similar distance drifts
heavily. Robustness, where it exists, looks like a property of \emph{what kind of knowledge}
a response draws on, not how far the prompt sits from the training domain.

\para{Generic fine-tuning drift dominates adversarial intent at small scale.} The \insecure
and \secure runs drifted identically (global $+0.004$). Emergent misalignment is, per
\citet{betley2025em}, a property that appears strongly only in larger models; a 135M model
under 30 LoRA steps shows only the stylistic adaptation common to \emph{any} fine-tune. The
methodological upshot is the paper's most transportable contribution: \emph{trajectory
divergence must be differenced against a matched benign-fine-tuning control to isolate
adversarial effects}. Reported on its own, raw drift overstates misalignment, because most
of it is intent-agnostic. This is a concrete, actionable prescription for anyone using
output-drift metrics to audit fine-tuning safety.

\para{How this answers the research question.} At the tested scale, language models do
\emph{not} show distance-graded robustness. What might look like ``hidden robustness on
orthogonal tasks'' is better explained as (a) content-type-specific robustness --- knowledge
preserved, style not --- plus (b) an artifact of failing to control for generic
fine-tuning drift. The trajectory-divergence and reversal-rate apparatus works and is
informative; its \emph{uncontrolled} output is simply not a misalignment signal.

\subsection{Limitations}
\label{sec:limitations}

We are explicit about scope: this is a small-scale CPU proof-of-concept and our conclusions
are about \emph{the method and its behavior at small scale}, not definitive claims about
frontier models.

\begin{itemize}[leftmargin=*,itemsep=1pt,topsep=2pt]
    \item \textbf{Severe scale reduction.} We used a 135M model (not a 7B model), LoRA
    $r=16$ (not full fine-tuning), 30 steps, 15 probes, and 3 checkpoints, all forced by a
    CPU-only, $\sim$2\,GB-RAM budget; full-fine-tuning attempts were out-of-memory-killed.
    \item \textbf{Low statistical power.} With $n=15$ (3 per tier), the Spearman CI spans
    $[-0.61, 0.51]$. We cannot detect a moderate true correlation; our null result is
    \emph{consistent with} but not \emph{proof of} no effect.
    \item \textbf{Emergent misalignment likely needs scale.} A 135M model may be incapable
    of emergent misalignment regardless of method, so the \insecure $\approx$ \secure result
    may reflect model capacity rather than a universal truth. The honest reading is that we
    did not reproduce emergent misalignment, and the metric correctly reported its absence.
    \item \textbf{Single seed, greedy decoding.} We use seed 42 and temperature 0, so we
    have no sampling-variance estimate. Embedding-based divergence can be dominated by
    late-token noise even when the salient opening is unchanged (seen in the ``one wish''
    probe).
    \item \textbf{Distance proxy.} Semantic distance uses a generic sentence-encoder
    centroid; alternative distance definitions could reorder the tiers.
\end{itemize}

\subsection{Broader implications}
\label{sec:implications}

For safety auditing, the practical message is to treat raw fine-tuning drift with caution
and always pair it with a benign-fine-tuning control before attributing drift to adversarial
intent. For robustness research, our results redirect attention from \emph{task distance} to
\emph{content type}: factual and technical knowledge appears to be a more reliable robustness
axis than orthogonality, and is worth studying directly. The reversal-rate signal, though
small here, suggests that drift dynamics --- overshoot and partial reversion --- carry
information beyond endpoint magnitude, supporting the trajectory view of fine-tuning safety.
